\subsection*{先有体积测量，再寻找代表它的箭头}
固定底面边 $u,v$，让第三条边 $w$ 变化。三边围成的有向体积是 $w$ 的函数。把 $w$ 拆成两段，两份有向体积相加；把它放大，体积同比例变化。因此体积是对 $w$ 的线性测量，恰好落在点积与对偶的框架中。

标准三维内积中，每个线性测量都能写成某个固定向量 $n$ 与 $w$ 的点积。怎样找 $n$？分别输入三个标准基，读数就是它的三个坐标。展开这三个行列式，得到的正是叉积公式；复杂的分量来自一个简单的代表问题。

\subsection*{垂直与长度从同一要求产生}
若 $w$ 落在底面里，三条边共面，体积为零。所以 $n\cdot w=0$ 对底面内所有向量成立，迫使 $n$ 垂直底面。若 $w$ 是沿正法线的单位向量，体积等于底面积，所以 $n$ 的长度也必须等于底面积。几何定义并非另加的装饰。

例如底边为 $(1,0,0),(0,2,0)$，第三边为 $(a,b,3)$。改变 $a,b$ 只使顶面沿底面剪切，体积始终为 $6$。代表测量的 $n=(0,0,2)$ 点乘第三边，直接读出同一结果。

值得保留的方法是：遇到一个对某变量线性的复杂运算，先问它对基向量分别做什么。矩阵列、点积的测量向量、叉积法向量，都是用少量基上的信息描述全部输入。学会这条思路，比记住三张看似独立的公式表更能帮助迁移。
